%% CHAPTER 12 DEFINITE INTEGRALS -- dictated by Maximus (25 Oct 2026); diagrams adapted from h2-math-notes/aoi/notes (figstyles sketch/curve/solX/solY)
\providecommand{\cautionbox}[2][CAUTION]{\begin{info}[title={#1}]#2\end{info}}
\providecommand{\ddx}{\,\mathrm{d}x}\providecommand{\ddy}{\,\mathrm{d}y}
\pgfplotsset{dimini/.style={sketch,width=3.7cm,height=3.1cm,every node/.append style={font=\scriptsize}}}
\newcommand{\rsum}[1]{% #1 = 0 left, 1 right
\begin{tikzpicture}[x=0.95cm,y=0.75cm,font=\scriptsize]
\def\f(##1){0.5+0.18*(##1)*(##1)}
\foreach \k in {0,1,2}{\pgfmathsetmacro\xl{0.8+0.5*\k}\pgfmathsetmacro\xr{\xl+0.5}\pgfmathsetmacro\xs{#1==0?\xl:\xr}
 \fill[blue!22,draw=blue!60!black,thin] (\xl,0) rectangle (\xr,{0.5+0.18*\xs*\xs});}
\foreach \k in {0,1}{\pgfmathsetmacro\xl{3.1+0.5*\k}\pgfmathsetmacro\xr{\xl+0.5}\pgfmathsetmacro\xs{#1==0?\xl:\xr}
 \fill[blue!22,draw=blue!60!black,thin] (\xl,0) rectangle (\xr,{0.5+0.18*\xs*\xs});}
\node at (2.7,0.55) {$\cdots$};
\draw[->,thick] (-0.2,0)--(4.4,0) node[right]{$x$}; \draw[->,thick] (0,-0.2)--(0,4.0) node[above]{$y$};
\draw[blue,semithick,domain=0.5:4.3,samples=60] plot (\x,{0.5+0.18*\x*\x}) node[right]{$y=\mathrm{f}(x)$};
\node[below] at (0.8,0){$a$}; \node[below] at (1.3,0){$a\!+\!h$}; \node[below] at (1.8,-0.3){$a\!+\!2h$};
\node[below] at (2.7,0){$\cdots$}; \node[below] at (4.1,0){$b$};
\end{tikzpicture}}
\newcommand{\vsolo}[1]{\adjustbox{totalheight=1.9cm,max width=6.6cm}{#1}}\newcommand{\vpair}[2]{\adjustbox{totalheight=1.9cm,max width=6.6cm}{#1\;\raisebox{0.15ex}{\Large$-$}\;#2}}
\newcommand{\cellfig}[2]{\begin{minipage}[t]{\linewidth}\centering\vspace{0pt}#1\\[-0.5mm]{\footnotesize$\displaystyle #2$}\vspace{1mm}\end{minipage}}
\clearpage
\inchap{Definite Integrals}
\cdisp{\int_{a}^{b}\mathrm{f}(x)\ddx=\Big[\mathrm{F}(x)\Big]_{a}^{b}=\mathrm{F}(b)-\mathrm{F}(a)}
\section{Area as a Limit of a Sum}
{\renewcommand{\arraystretch}{1.3}\centering\begin{tabular}{|>{\centering\arraybackslash}m{7.6cm}|>{\centering\arraybackslash}m{7.6cm}|}\hline
\textbf{Left-hand sum}&\textbf{Right-hand sum}\\\hline
\rsum{0}&\rsum{1}\\\hline
$\displaystyle\lim_{n\to\infty}\sum_{r=0}^{n-1}\frac{b-a}{n}\,\mathrm{f}\left(a+\frac{r(b-a)}{n}\right)=\int_{a}^{b}\mathrm{f}(x)\ddx$&$\displaystyle\lim_{n\to\infty}\sum_{r=1}^{n}\frac{b-a}{n}\,\mathrm{f}\left(a+\frac{r(b-a)}{n}\right)=\int_{a}^{b}\mathrm{f}(x)\ddx$\\[3mm]\hline
Increasing $\mathrm{f}(x)\Rightarrow$ underestimate\newline Decreasing $\mathrm{f}(x)\Rightarrow$ overestimate&Increasing $\mathrm{f}(x)\Rightarrow$ overestimate\newline Decreasing $\mathrm{f}(x)\Rightarrow$ underestimate\\\hline
\end{tabular}\par}
\medskip
%% ---------------- dx / dy table
\newcommand{\gX}[2]{\begin{tikzpicture}\begin{axis}[dimini,xmin=-0.3,xmax=4.3,#1]#2\end{axis}\end{tikzpicture}}
{\renewcommand{\arraystretch}{1.15}\setlength{\tabcolsep}{3pt}\centering\begin{longtable}{|>{\raggedright\arraybackslash}m{2.3cm}|>{\centering\arraybackslash}m{7.15cm}|>{\centering\arraybackslash}m{7.15cm}|}\hline
\textbf{Index}&$\boldsymbol{\mathrm{d}x}$&$\boldsymbol{\mathrm{d}y}$\\\hline\endhead
Area under curve&
\cellfig{\gX{ymin=-0.3,ymax=3}{\fill[blue!22] (axis cs:1,0)--plot[domain=1:3.4,samples=30](axis cs:\x,{1+0.15*\x*\x})--(axis cs:3.4,0)--cycle;\addplot[curve,domain=0.4:3.9]{1+0.15*x*x};\node[below] at (axis cs:1,0){$a$};\node[below] at (axis cs:3.4,0){$b$};}}{\int_{a}^{b}y\ddx}&
\cellfig{\gX{ymin=-0.3,ymax=3.4}{\fill[blue!22] (axis cs:0,0.8)--plot[domain=0.8:2.8,samples=30](axis cs:{1+0.25*\x*\x},\x)--(axis cs:0,2.8)--cycle;\addplot[curve,variable=t,domain=0.4:3.2]({1+0.25*t*t},{t});\node[left] at (axis cs:0,0.8){$c$};\node[left] at (axis cs:0,2.8){$d$};}}{\int_{c}^{d}x\ddy}\\\hline
Negative areas&
\cellfig{\gX{ymin=-3,ymax=0.5}{\fill[blue!22] (axis cs:1,0)--plot[domain=1:3.4,samples=30](axis cs:\x,{-1-0.15*\x*\x})--(axis cs:3.4,0)--cycle;\addplot[curve,domain=0.4:3.9]{-1-0.15*x*x};\node[above] at (axis cs:1,0){$a$};\node[above] at (axis cs:3.4,0){$b$};}}{-\int_{a}^{b}y\ddx}&
\cellfig{\gX{xmin=-4.3,xmax=0.4,ymin=-0.3,ymax=3.4}{\fill[blue!22] (axis cs:0,0.8)--plot[domain=0.8:2.8,samples=30](axis cs:{-1-0.25*\x*\x},\x)--(axis cs:0,2.8)--cycle;\addplot[curve,variable=t,domain=0.4:3.2]({-1-0.25*t*t},{t});\node[right] at (axis cs:0,0.8){$c$};\node[right] at (axis cs:0,2.8){$d$};}}{-\int_{c}^{d}x\ddy}\\\hline
Areas between two graphs&
\cellfig{\gX{ymin=-0.3,ymax=3.4,width=3.1cm}{\fill[blue!22] plot[domain=1:3.4,samples=30](axis cs:\x,{2.2+0.1*\x})--plot[domain=3.4:1,samples=30](axis cs:\x,{0.4+0.12*(\x-2)^2})--cycle;\addplot[curve,domain=0.5:3.9]{2.2+0.1*x};\addplot[line2,red!75!black,domain=0.5:3.9,samples=40]{0.4+0.12*(x-2)^2};\draw[dashed,thin](axis cs:1,0)--(axis cs:1,2.3);\draw[dashed,thin](axis cs:3.4,0)--(axis cs:3.4,2.54);\node[below] at (axis cs:1,0){$a$};\node[below] at (axis cs:3.4,0){$b$};}\,%
\gX{ymin=-1.8,ymax=2,width=3.1cm}{\fill[blue!22] plot[domain=1:3.4,samples=30](axis cs:\x,{1.2+0.1*\x})--plot[domain=3.4:1,samples=30](axis cs:\x,{-1.2+0.15*(\x-2)^2})--cycle;\addplot[curve,domain=0.5:3.9]{1.2+0.1*x};\addplot[red!75!black,semithick,domain=0.5:3.9,samples=40]{-1.2+0.15*(x-2)^2};\node[above] at (axis cs:1,0){$a$};\node[above] at (axis cs:3.6,0){$b$};}\,%
\gX{ymin=-3.4,ymax=0.4,width=3.1cm}{\fill[blue!22] plot[domain=1:3.4,samples=30](axis cs:\x,{-0.5-0.1*\x})--plot[domain=3.4:1,samples=30](axis cs:\x,{-2.6+0.12*(\x-2)^2})--cycle;\addplot[curve,domain=0.5:3.9]{-0.5-0.1*x};\addplot[red!75!black,semithick,domain=0.5:3.9,samples=40]{-2.6+0.12*(x-2)^2};\node[above] at (axis cs:1,0){$a$};\node[above] at (axis cs:3.4,0){$b$};}}%
{\text{All 3 cases: }\int_{a}^{b}\left(y_{1}-y_{2}\right)\ddx}&
\cellfig{\gX{xmin=-0.3,xmax=3.4,ymin=-0.3,ymax=4.3,width=3.1cm}{\fill[blue!22] plot[domain=1:3.4,samples=30](axis cs:{2.2+0.1*\x},\x)--plot[domain=3.4:1,samples=30](axis cs:{0.4+0.12*(\x-2)^2},\x)--cycle;\addplot[curve,variable=t,domain=0.5:3.9]({2.2+0.1*t},{t});\addplot[red!75!black,semithick,variable=t,domain=0.5:3.9,samples=40]({0.4+0.12*(t-2)^2},{t});\node[left] at (axis cs:0,1){$c$};\node[left] at (axis cs:0,3.4){$d$};}\,%
\gX{xmin=-1.8,xmax=2,ymin=-0.3,ymax=4.3,width=3.1cm}{\fill[blue!22] plot[domain=1:3.4,samples=30](axis cs:{1.2+0.1*\x},\x)--plot[domain=3.4:1,samples=30](axis cs:{-1.2+0.15*(\x-2)^2},\x)--cycle;\addplot[curve,variable=t,domain=0.5:3.9]({1.2+0.1*t},{t});\addplot[red!75!black,semithick,variable=t,domain=0.5:3.9,samples=40]({-1.2+0.15*(t-2)^2},{t});\node[right] at (axis cs:0,1){$c$};\node[right] at (axis cs:0,3.6){$d$};}\,%
\gX{xmin=-3.4,xmax=0.4,ymin=-0.3,ymax=4.3,width=3.1cm}{\fill[blue!22] plot[domain=1:3.4,samples=30](axis cs:{-0.5-0.1*\x},\x)--plot[domain=3.4:1,samples=30](axis cs:{-2.6+0.12*(\x-2)^2},\x)--cycle;\addplot[curve,variable=t,domain=0.5:3.9]({-0.5-0.1*t},{t});\addplot[red!75!black,semithick,variable=t,domain=0.5:3.9,samples=40]({-2.6+0.12*(t-2)^2},{t});\node[right] at (axis cs:0,1){$c$};\node[right] at (axis cs:0,3.4){$d$};}}%
{\text{All 3 cases: }\int_{c}^{d}\left(x_{1}-x_{2}\right)\ddy}\\
&{\scriptsize $y_1$ (blue) above $y_2$ (red)}&{\scriptsize $x_1$ (blue) right of $x_2$ (red)}\\\hline
Volume rotated about axis&
\cellfig{\gX{ymin=-0.3,ymax=3}{\fill[blue!22] (axis cs:1,0)--plot[domain=1:3.4,samples=30](axis cs:\x,{1+0.15*\x*\x})--(axis cs:3.4,0)--cycle;\addplot[curve,domain=0.4:3.9]{1+0.15*x*x};\node[below] at (axis cs:1,0){$a$};\node[below] at (axis cs:3.4,0){$b$};}\\[1mm]\vsolo{\solX[0.9]{1+0.15*x*x}{1}{3.4}{1.15}{2.73}{4}{3.2}{blue!70!black}}}{\text{About the $x$-axis: }\pi\int_{a}^{b}y^{2}\ddx}&
\cellfig{\gX{ymin=-0.3,ymax=3.4}{\fill[blue!22] (axis cs:0,0.8)--plot[domain=0.8:2.8,samples=30](axis cs:{1+0.25*\x*\x},\x)--(axis cs:0,2.8)--cycle;\addplot[curve,variable=t,domain=0.4:3.2]({1+0.25*t*t},{t});\node[left] at (axis cs:0,0.8){$c$};\node[left] at (axis cs:0,2.8){$d$};}\\[1mm]\vsolo{\solY{1+0.25*x*x}{0.8}{2.8}{1.16}{2.96}{3.4}{3.4}{blue!70!black}}}{\text{About the $y$-axis: }\pi\int_{c}^{d}x^{2}\ddy}\\\hline
Volume generated by region between two curves&
\cellfig{\gX{ymin=-0.3,ymax=3.4}{\fill[blue!22] plot[domain=1:3.4,samples=30](axis cs:\x,{2.2+0.1*\x})--plot[domain=3.4:1,samples=30](axis cs:\x,{0.4+0.12*(\x-2)^2})--cycle;\addplot[curve,domain=0.5:3.9]{2.2+0.1*x};\addplot[red!75!black,semithick,domain=0.5:3.9,samples=40]{0.4+0.12*(x-2)^2};\node[below] at (axis cs:1,0){$a$};\node[below] at (axis cs:3.4,0){$b$};}\\[1mm]\vpair{\solX[0.9]{2.2+0.1*x}{1}{3.4}{2.3}{2.54}{4}{3}{blue!70!black}}{\solX[0.9]{0.4+0.12*(x-2)^2}{1}{3.4}{0.52}{0.63}{4}{3}{red!75!black}}}{\pi\int_{a}^{b}\left(y_{1}^{2}-y_{2}^{2}\right)\ddx}&
\cellfig{\gX{xmin=-0.3,xmax=3.4,ymin=-0.3,ymax=4.3}{\fill[blue!22] plot[domain=1:3.4,samples=30](axis cs:{2.2+0.1*\x},\x)--plot[domain=3.4:1,samples=30](axis cs:{0.4+0.12*(\x-2)^2},\x)--cycle;\addplot[curve,variable=t,domain=0.5:3.9]({2.2+0.1*t},{t});\addplot[red!75!black,semithick,variable=t,domain=0.5:3.9,samples=40]({0.4+0.12*(t-2)^2},{t});\node[left] at (axis cs:0,1){$c$};\node[left] at (axis cs:0,3.4){$d$};}\\[1mm]\vpair{\solY{2.2+0.1*x}{1}{3.4}{2.3}{2.54}{3.2}{4}{blue!70!black}}{\solY{0.4+0.12*(x-2)^2}{1}{3.4}{0.52}{0.63}{3.2}{4}{red!75!black}}}{\pi\int_{c}^{d}\left(x_{1}^{2}-x_{2}^{2}\right)\ddy}\\\hline
\end{longtable}\par}
When finding \textbf{area}, we generally have the choice to integrate via $\mathrm{d}x$ or $\mathrm{d}y$. Some considerations:
\begin{itemize}
\item $\mathrm{d}x$ or $\mathrm{d}y$ easier to integrate
\item do we need to add or subtract regular shapes like rectangles
\item is it even possible to make $x$ (or $y$) the subject?
\end{itemize}
For \textbf{volume}, we cannot choose which variable to integrate with respect to.
\begin{itemize}
\item Region rotated about the $x$-axis: $\displaystyle\pi\int_{a}^{b}y^{2}\ddx$
\item Region rotated about the $y$-axis: $\displaystyle\pi\int_{c}^{d}x^{2}\ddy$
\end{itemize}
\begin{info}[title={Important Note: dealing with conics}]
When making $y$ the subject, the sign of the square root indicates which half of the conic is being considered.\par\smallskip
\textbf{Example}\quad $(x-1)^{2}+(y-2)^{2}=9$\par
{\renewcommand{\arraystretch}{1.1}\centering\begin{tabular}{>{\centering\arraybackslash}m{5cm}>{\centering\arraybackslash}m{3.2cm}}
$\displaystyle y=2+\sqrt{9-(x-1)^{2}}$&\begin{tikzpicture}[scale=0.32]\draw[->](-2.6,0)--(4.8,0) node[right]{\tiny$x$};\draw[->](0,-1.6)--(0,5.8) node[above]{\tiny$y$};\draw[gray!60,dashed](1,2) circle(3);\draw[blue,ultra thick](1,2)++(0:3) arc(0:180:3);\fill(1,2) circle(2pt);\end{tikzpicture}\\
$\displaystyle y=2-\sqrt{9-(x-1)^{2}}$&\begin{tikzpicture}[scale=0.32]\draw[->](-2.6,0)--(4.8,0) node[right]{\tiny$x$};\draw[->](0,-1.6)--(0,5.8) node[above]{\tiny$y$};\draw[gray!60,dashed](1,2) circle(3);\draw[blue,ultra thick](1,2)++(180:3) arc(180:360:3);\fill(1,2) circle(2pt);\end{tikzpicture}\\
$\displaystyle x=1+\sqrt{9-(y-2)^{2}}$&\begin{tikzpicture}[scale=0.32]\draw[->](-2.6,0)--(4.8,0) node[right]{\tiny$x$};\draw[->](0,-1.6)--(0,5.8) node[above]{\tiny$y$};\draw[gray!60,dashed](1,2) circle(3);\draw[blue,ultra thick](1,2)++(-90:3) arc(-90:90:3);\fill(1,2) circle(2pt);\end{tikzpicture}\\
$\displaystyle x=1-\sqrt{9-(y-2)^{2}}$&\begin{tikzpicture}[scale=0.32]\draw[->](-2.6,0)--(4.8,0) node[right]{\tiny$x$};\draw[->](0,-1.6)--(0,5.8) node[above]{\tiny$y$};\draw[gray!60,dashed](1,2) circle(3);\draw[blue,ultra thick](1,2)++(90:3) arc(90:270:3);\fill(1,2) circle(2pt);\end{tikzpicture}
\end{tabular}\par}
Note: this concept can be used for any conic. That is, circles, ellipses, hyperbolas, parabolas.
\end{info}
\textbf{Notes on volume}
\begin{itemize}
\item ``Rotate $\pi$ radians about \dots'' does not usually mean the formula has to be divided by 2. It just means that rotating only $180^{\circ}$ is sufficient to generate the full solid. The volume remains as $\pi\int_a^b y^2\ddx$.
\item To find the volume of rotating regions ``below'' the axis, there is no need to put a negative sign in front:
\cdisp{\textcolor{red}{\text{\ding{55}}}\quad-\pi\int_{a}^{b}y^{2}\ddx\qquad\qquad\textcolor{green!45!black}{\text{\ding{51}}}\quad\pi\int_{a}^{b}y^{2}\ddx}
\item Common mistake for volume between two graphs:
\cdisp{\textcolor{red}{\text{\ding{55}}}\quad\pi\int_{a}^{b}\left(y_{1}-y_{2}\right)^{2}\ddx\qquad\qquad\textcolor{green!45!black}{\text{\ding{51}}}\quad\pi\int_{a}^{b}\left(y_{1}^{2}-y_{2}^{2}\right)\ddx}
\end{itemize}
\begin{info}[title=Important Note]
Take note if the question asks for exact or not. If they did not, please evaluate the integral using GC rather than integration by hand.
\end{info}
\begin{info}[title=Remark]
If an integral seems very difficult/impossible to integrate, then check if the question wants exact form OR check if you have integrated something similar in earlier parts.
\end{info}
\subsection*{\fhl{12.1.1}\quad Regular Shapes in Area and Volume}\addcontentsline{toc}{subsection}{\protect\numberline{12.1.1}Regular Shapes in Area and Volume}
Sometimes we may encounter familiar shapes like squares, rectangles, triangles, trapeziums, cones, and cylinders. It will be useful to use the formulas for these shapes/solids for efficiency.
\begin{multicols}{3}
\centering
\begin{tikzpicture}[scale=0.8]\draw[thick,fill=blue!15](0,0) rectangle (1.8,1.8);\node[below] at (0.9,0){$a$};\node[left] at (0,0.9){$a$};\end{tikzpicture}\\Square: $\text{Area}=a^{2}$\par\medskip
\begin{tikzpicture}[scale=0.8]\draw[thick,fill=blue!15](0,0) rectangle (2.8,1.6);\node[below] at (1.4,0){$l$};\node[left] at (0,0.8){$b$};\end{tikzpicture}\\Rectangle: $\text{Area}=lb$\par\medskip
\columnbreak
\begin{tikzpicture}[scale=0.8]\draw[thick,fill=blue!15](0,0)--(2.8,0)--(1.0,1.8)--cycle;\draw[dashed](1.0,1.8)--(1.0,0);\draw(1.0,0.2)--(1.2,0.2)--(1.2,0);\node[below] at (1.4,0){$b$};\node[right] at (1.0,0.9){$h$};\end{tikzpicture}\\Triangle: $\text{Area}=\frac12bh$\par\medskip
\begin{tikzpicture}[scale=0.8]\draw[thick,fill=blue!15](0,0)--(3,0)--(2.3,1.6)--(0.6,1.6)--cycle;\draw[dashed](0.6,1.6)--(0.6,0);\node[below] at (1.5,0){$b$};\node[above] at (1.45,1.6){$a$};\node[left] at (0.6,0.8){$h$};\end{tikzpicture}\\Trapezium: $\text{Area}=\frac12(a+b)h$\par\medskip
\columnbreak
\begin{tikzpicture}[scale=0.8]\fill[blue!15](-1,0)--(0,2.2)--(1,0) arc(0:-180:1 and 0.3);\draw[thick](-1,0)--(0,2.2)--(1,0);\draw[thick](1,0) arc(0:-180:1 and 0.3);\draw[dashed](1,0) arc(0:180:1 and 0.3);\draw[dashed](0,0)--(0,2.2);\draw[dashed](0,0)--(1,0);\node[right] at (0,1.1){$h$};\node[above] at (0.5,0){$r$};\end{tikzpicture}\\Cone: $\text{Volume}=\frac13\pi r^{2}h$\par\medskip
\begin{tikzpicture}[scale=0.8]\fill[blue!15](-1,0)--(-1,1.8) arc(180:360:1 and 0.3)--(1,0) arc(0:-180:1 and 0.3);\draw[thick,fill=blue!8](0,1.8) ellipse(1 and 0.3);\draw[thick](-1,0)--(-1,1.8);\draw[thick](1,0)--(1,1.8);\draw[thick](1,0) arc(0:-180:1 and 0.3);\draw[dashed](1,0) arc(0:180:1 and 0.3);\draw[dashed](0,1.8)--(1,1.8);\node[above] at (0.5,1.8){$r$};\node[right] at (1,0.9){$h$};\end{tikzpicture}\\Cylinder: $\text{Volume}=\pi r^{2}h$
\end{multicols}
\subsection*{\fhl{12.1.2}\quad Integration with Modulus}\addcontentsline{toc}{subsection}{\protect\numberline{12.1.2}Integration with Modulus}
To integrate with modulus, we determine where the function inside the modulus is positive and where it is negative, then split it up.
\exlab
\cdisp{\int_{0}^{\pi}\left|x\sin x\cos x\right|\ddx&=\int_{0}^{\frac{\pi}{2}}{\color{blue}(}x\sin x\cos x{\color{blue})}\ddx+\int_{\frac{\pi}{2}}^{\pi}{\color{red}-(}x\sin x\cos x{\color{red})}\ddx\\
&=\int_{0}^{\frac{\pi}{2}}\frac12x\sin2x\ddx+\int_{\frac{\pi}{2}}^{\pi}{\color{red}-}\frac12x\sin2x\ddx\\
&=\ldots}
